\section{Peripheral obstructions with chordal maximal-cube graphs}
\label{sec:rper-chordal-sun-obstructions}

The maximal-cube graph does not determine membership in $\Rper$, even
among ample classes for which that graph is chordal. We construct an
infinite forbidden family with this property. Its members nevertheless
have explicit corner peelings and acyclic representation maps. Thus
chordal maximal-cube graphs do not reduce the peripheral obstruction
problem to corner-peeling obstructions.

Here the graph has maximal positioned cubes as vertices, with edges for
nonempty intersection. It differs from the shattered-pair graph, whose
vertices are coordinates. The ACC criterion in the
\href{../../../YangDim/yang_complex_dictionary.tex}{Yang-complex dictionary}
requires medianity and chordality of that latter graph. Directly, ACC
requires a join tree of maximal cubes, with connected occurrence nodes for
each concept. A chordal intersection graph alone does not ensure such a
tree: the ample downset $\{0,1\}^3\setminus\{111\}$ has three maximal
squares with intersection graph $K_3$. Its three concepts of weight one
each force a different pair of squares to be adjacent in a join tree,
which would require all three edges. It is therefore not ACC.

Fix $m\ge4$, with indices taken modulo $m$. The coordinate set is
\[
 E_m=\{(i,j):j=i-1\text{ or }i+1\},\qquad |E_m|=2m.
\]
We identify a word with its set of coordinates having value one. Define
its two incoming and outgoing coordinate pairs by
\[
 A_i=\{(i-1,i),(i+1,i)\},\qquad
 B_i=\{(i,i-1),(i,i+1)\},
\]
and put $S_i=A_i\cup A_{i+1}$ and $W_i=A_i\cup B_i$.
For $J\subseteq\mathbb Z_m$, let
\begin{equation}
 D_m=\bigcup_{i\in\mathbb Z_m}2^{S_i},\qquad
 H_{m,J}=D_m\cup\{W_i:i\in J\},\qquad H_m=H_{m,\mathbb Z_m}.
 \label{eq:rper-sun-family}
\end{equation}
The indices in $E_m$ are coordinate names, not edges of the concept graph.
The set $D_m$ is a Boolean downset and therefore belongs to $\Rper$.
Each added word completes a square based at $A_i$, with free directions
$B_i$.

\begin{theorem}[An infinite chordal forbidden family]
\label{thm:rper-chordal-sun-obstructions}
For every $m\ge4$, the class $H_m$ is an ample forbidden pc-minor of
$\Rper$, of isometric dimension $2m$, VC dimension four and order $13m+1$.
Every proper subfamily $H_{m,J}$ with $J\subsetneq\mathbb Z_m$ belongs
to $\Rper$.

The complete maximal cubes of $H_m$ are
\[
 C_i=2^{S_i},\qquad
 L_i=\{A_i\cup T:T\subseteq B_i\}\qquad(i\in\mathbb Z_m).
\]
The $C_i$ form a clique in their intersection graph. The $L_i$ are
pairwise disjoint, and $L_i$ meets exactly $C_{i-1},C_i$.
Thus this graph is an $m$-sun and is chordal: delete the simplicial
vertices $L_i$, then the remaining clique.

The graph of $H_m$ has no cut vertex. Deleting the words $W_i$ in any
order, followed by the words of $D_m$ in decreasing cardinality, is a
corner peeling. Its explicit acyclic ordinary representation is
\begin{equation}
 r(U)=U\quad(U\in D_m),\qquad r(W_i)=B_i.
 \label{eq:rper-sun-representation}
\end{equation}
For every coordinate the two restrictions have orders $16$ and $13m-15$,
and the contraction has order $13m-14$.
\end{theorem}

\paragraph{Proof strategy.}
The downset has $12m+1$ words. Each $W_i$ adds precisely the shattered
support $B_i$, and is a degree-two corner. Peripheral recursion fails
only when all the extra words are present. If one is missing, choose a
neighboring present index and use the coordinate directed from it toward
the missing index. One section deletes that present word; the other is a
downset with at most one pendant word. This argument works through the
successive zero restrictions and supplies positivity of every proper
subfamily. A contraction absorbs one of the added words into the downset,
so it is the contraction of a proper positive subfamily.

\begin{proof}
\emph{The downset and the new supports.}
The pairs $A_i$ partition $E_m$. Each $S_i$ consists of two consecutive
pairs. The downset $D_m$ has $2m$ singletons, $5m$ pairs, $4m$ triples
and $m$ four-sets, as well as the empty set. Indeed each $S_i$ contributes
six pairs, and its one pair $A_i$ is also in $S_{i-1}$; these are the
only repetitions of pairs between different cores. No triple or four-set
is repeated. Hence $|D_m|=12m+1$.

The two coordinates of $B_i$ have targets $i-1$ and $i+1$, which are not
consecutive for $m\ge4$. Thus $B_i\notin D_m$ and $W_i\notin D_m$.
For either $b\in B_i$, however, $W_i\setminus\{b\}$ lies in one of the
cores $2^{S_{i-1}},2^{S_i}$. Removing an incoming coordinate from $W_i$
leaves both coordinates of $B_i$, so does not give a word of $D_m$.
Different $W_i$ have the same cardinality and are never neighbors.
It follows that $W_i$ is a corner of degree two, with incident square
$L_i$, in every $H_{m,J}$ containing it.

For completeness, the exact shattering identity is
\begin{equation}
 \strc(H_{m,J})=D_m\ \dot\cup\ \{B_i:i\in J\}.
 \label{eq:rper-sun-shadow}
\end{equation}
Consider a newly shattered support $T$. If every new word has its
$T$-trace already present in $D_m$, no new trace appears. Otherwise some
$W_i|_T$ was absent. Since removing either coordinate of $B_i$ from
$W_i$ gives an old word, necessarily $B_i\subseteq T$.
Only $W_i$ among all the words of $H_{m,J}$ contains both coordinates
of $B_i$: for $j\ne i$, $W_j$ contains at most one of the two outgoing
coordinates with source $i$. Consequently any additional coordinate
of $T\setminus B_i$ has a forced value in the trace with $B_i=11$,
so $T$ cannot be shattered. The new support must therefore be $B_i$,
and its full square $L_i$ is present. All supports in $D_m$ have their
cubes based at zero. This proves \eqref{eq:rper-sun-shadow} and ampleness,
as well as $|H_{m,J}|=12m+1+|J|$ and VC dimension four.
Every coordinate is active since all singletons belong to $D_m$.

The maximal shattered supports in \eqref{eq:rper-sun-shadow} are exactly
$S_i$ and the selected $B_i$. In an ample class, maximal cubes are the
unique cubes on maximal shattered supports. To recall why, the anchors
of cubes on a fixed support form an ample full reduction, hence are
connected. A maximal cube's anchor cannot have a neighbor in that
reduction, which would extend the cube. If its support were not maximal shattered, a larger strongly shattered
cube would give at least two anchors in that full reduction, contradicting
connectedness and the isolated anchor. Thus the support is maximal, with
only one position; the converse follows immediately.
This identifies the complete cube list as claimed.
Every $C_i$ contains zero. A word of $L_i$ contains $A_i$, so $L_i$ can
meet only $C_{i-1}$ and $C_i$. It meets both at the word $A_i$.
For $j\ne i$, the set $A_j$ is not contained in $W_i$: at most one of
its two coordinates has source $i$. Thus $L_i\cap L_j=\varnothing$.
These observations give the asserted chordal intersection graph.
Each core four-cube remains connected after deletion of any one vertex,
and consecutive core intersections are four-vertex squares, which remain
nonempty. Thus $D_m$ has no cut vertex. Every additional word has two
neighbors in $D_m$, so adding them preserves the absence of cut vertices.
The cubes just listed equivalently describe the complete extentures of
$H_m^c$, not those of $H_m$: the core cubes fix zero outside $S_i$;
the leaf cubes fix one on $A_i$ and zero outside $A_i\cup B_i$.

\emph{A proper selected family has a peripheral recursion.}
We prove the slightly stronger statement needed to iterate restrictions.
Let $P\subseteq E_m$ be coordinates fixed to zero. The corresponding
section of $H_{m,J}$ consists of the restricted downset and those $W_i$
with $i\in J$ and $W_i\cap P=\varnothing$. Call these indices live.
If the live set is proper, this section belongs to $\Rper$.
Induct on the number of live indices. With none, the section is a downset.
Otherwise choose adjacent indices $f,e$ with $f$ live and $e$ not live;
this is possible because the cycle is connected. The coordinate
$c=(f,e)$ is not in $P$, since $c\in W_f$.
Among the added words, only $W_f,W_e$ contain $c$, and $W_e$ is absent.
Lowering $c$ in $W_f$ gives a word of the downset. Therefore this
coordinate is downward peripheral.

In its zero section, $f$ ceases to be live, so induction applies.
In its one section, the downset portion comes only from the two cores
containing $A_e$, with $P\cup\{c\}$ removed; it is a downset after
$c$ is suppressed. Its only additional word is $W_f\setminus\{c\}$.
Write $h$ for the neighbor of $f$ other than $e$ and $t=(f,h)$.
The coordinate $t$ occurs in neither of those two core supports: $h$ is
neither $e$ nor a neighbor of $e$, since $m\ge4$.
The extra word is $A_f\cup\{t\}$ and its neighbor $A_f$ lies in the
core shared by $e,f$. None of these coordinates has been fixed to zero,
because $f$ is live. Hence the one section is a downset with a single
pendant word, attached in the previously constant coordinate $t$.
It belongs to $\Rper$ by a singleton peripheral expansion of the
downset. Both children are positive, completing the induction.
Taking $P=\varnothing$ proves $H_{m,J}\in\Rper$ whenever $J$ is proper.

\emph{Failure of peripherality and positivity of every proper minor.}
For $c=(f,e)$, the leaf cube $L_e$ fixes $c=1$. Any core whose two target
indices omit $e$ fixes $c=0$. These are maximal cubes.
If the two coordinate sections were nested, a maximal cube fixed on
the smaller side could be duplicated across that coordinate, contradicting
maximality. Thus both fixed polarities exclude peripherality.
There is such a core for every $m\ge4$, so $H_m$ has no peripheral
coordinate and cannot belong to $\Rper$.

For either prescribed bit of $c$, only a proper subset of the added
words survives: exactly $W_f,W_e$ survive on the one side, and all the
other $W_i$ survive on the zero side. Each restriction is therefore the
same restriction of a proper $H_{m,J}$, and belongs to $\Rper$ by
\cref{lem:pcminor-closure}.
For contraction, $W_f\setminus\{c\}$ lies in the projected downset,
because $c\in B_f$. Hence
\[
 \pi_c H_m=\pi_c H_{m,\mathbb Z_m\setminus\{f\}}\in\Rper.
\]
All elementary proper images are positive, and their further pc-minors
remain positive by the same closure lemma. This proves forbiddenness.

The one section has two core three-cubes sharing a one-cube, hence
$8+8-2=14$ downset words, and the two distinct extra words give order16.
The zero section has order $13m+1-16=13m-15$.
Every downset word of the one section has its zero neighbor; among its
two extra words, only $W_f$ has that neighbor. Thus there are15 double
fibers, giving contraction order $13m+1-15=13m-14$.

Finally the degree-two corners can be deleted in any order, leaving the
downset. Removing downset words in decreasing cardinality finishes a
corner peeling. Its outmap is precisely \eqref{eq:rper-sun-representation}:
the added words point along $B_i$ into the downset, and the downset map
removes its positive coordinates. Its labels are the support family in
\eqref{eq:rper-sun-shadow}, each once. All edges point from one to zero,
so the map is acyclic and every contained cube has its unique zeroward
sink. Its outgoing cubes are $L_i$ and the lower cubes in $D_m$.
This proves the representation assertion.
\end{proof}

\begin{remark}[The same abstract cube graph with opposite membership]
\label{rem:rper-sun-graph-not-classifier}
There is also an $\Rper$ class with this same $m$-sun as its complete
maximal-cube graph. On coordinates $x_i,p_i$, take the downset generated
by the pairs $\{x_i,x_{i+1}\}$ and add the pendant words $\{x_i,p_i\}$.
The downset has $2m+1$ words and each pendant attachment is peripheral.
Its $m$ core squares form a clique, and its $m$ pendant edges are disjoint,
each meeting precisely its two neighboring core squares.
Thus it has $3m+1$ words, VC dimension two, and the same cube graph as
$H_m$, but belongs to $\Rper$.
The abstract intersection graph forgets the cube supports and positions.
The theorem is not an exhaustive characterization of all obstructions
with chordal maximal-cube graphs, nor a maximum-parent theorem for them.
The \href{../maxmin-forbidden/research/chordal_sun_parents.tex}
{maximum-generated companion} supplies convex realizations and actual
maximum parents for this entire particular family $H_m$, uniformly in $m$.
\end{remark}
